% Image credits for the photographs and AI illustrations of Book 4 (University, Year 2).
% Machine-readable license records live in images/book4/CREDITS.md; keep the
% two files in sync. The Book 4 writer owns this file: add one \item per
% photograph (inside the itemize, by chapter) as photographs land.
\chapter*{Image Credits}

The drawings, maps and charts in this book are original; maps and charts
drawn from published data name their source in the caption or in this
list. Photographs are used under their respective free licenses, with
thanks to their authors; full source links and license URLs for every
photograph are listed in \texttt{images/book4/CREDITS.md} in the book's
repository.

\bigskip
Alongside the photographs and the original drawings, some figures are
AI-generated illustrations, produced with OpenAI image generation from
prompts written by the book's editors, and reviewed for scientific
accuracy before inclusion. They show generic scenes only, never a real
place, a real person or a real specimen. The full prompt for every
generated image is recorded in \texttt{images/book4/ai/PROMPTS.md} in the
repository.

\bigskip
\noindent Photographs and data, by chapter:
\begin{itemize}\small
  \item Ch.~1 (the interstellar medium and star formation): the protoplanetary disc of HL Tauri, ALMA (ESO/NAOJ/NRAO), CC BY 4.0; the Pillars of Creation in the Eagle Nebula, NASA, ESA and the Hubble Heritage Team (STScI/AURA), public domain.
  \item Ch.~2 (planet formation and meteorites): the Allende carbonaceous chondrite and Widmanst\"atten pattern of the Gibeon iron meteorite, James St.~John, CC BY 2.0.
  \item Ch.~3 (elements in the Earth): lherzolite mantle xenolith in basalt, Kilbourne Hole, James St.~John, CC BY 2.0.
  \item Ch.~5 (magmas and their emplacement): dyke near Shiprock, Dexter Perkins, CC0 (Wikimedia Commons).
  \item Ch.~6 (volcanic eruptions): lava fountain on the Pu`u~`\=O`\=o cone of Kīlauea, J.\,D. Griggs, USGS, public domain; the eruption column of Mount St.~Helens, 18 May 1980, Austin Post, USGS, public domain; the eruption column of Mount Pinatubo, 12 June 1991, Dave Harlow, USGS, public domain.
  \item Ch.~7 (melting the mantle): dunite hand specimen, Mike Davis (GeoDIL), CC0.
  \item Ch.~8 (magmatic differentiation): chromitite layers of the Bushveld Complex, Kevin Walsh, CC BY 2.0; layered gabbros of the Skaergaard intrusion, C. Tegner, CC BY 4.0.
  \item Ch.~10 (metamorphic reactions and facies): lawsonite blueschist, Ward Creek, California, and eclogite, Almenningen, Norway, James St.~John, CC BY 2.0.
  \item Ch.~11 (thermobarometry and P--T--t paths): garnet schist, James St.~John, CC BY 2.0; coesite grain in eclogite, crossed polars, J.\,R. Smyth, CC BY-SA 1.0.
  \item Ch.~12 (crustal melting: migmatites and granites): migmatite outcrop in Nashua, New Hampshire, Kevin M. Gill, CC BY 4.0; leucogranite vein, Anne Burgess (Geograph), CC BY-SA 2.0.
  \item Ch.~13 (anatomy of a mountain belt: the Alps): pillow lavas of the Chenaillet ophiolite, Lionel Duchoiselle, CC BY-SA 3.0; the Matterhorn from Zermatt, Ximonic (Simo Räsänen), CC BY-SA 4.0.
  \item Ch.~14 (collision and Andean belts compared): the Tibetan Plateau between Nepal and Shigatse, Jan Kranendonk, CC BY-SA 4.0; the Bolivian Altiplano, Luca Galuzzi, CC BY-SA 2.5.
  \item Ch.~15 (ancient massifs and the Wilson cycle): Siccar Point, south-east cliff, dave souza, CC BY-SA 4.0.
  \item Ch.~16 (sedimentary basins compared): semi-submersible drilling rig in the North Sea, Erik Christensen, CC BY-SA 3.0.
  \item Ch.~17 (oceanic islands): Kure Atoll from the International Space Station, NASA, public domain.
  \item Ch.~18 (the Precambrian Earth): spinifex komatiite at Pyke Hill, Ontario, James St.~John, CC BY 2.0; cap carbonate of the Nuccaleena Formation, South Australia, James St.~John, CC BY 2.0.
  \item Ch.~19 (the Phanerozoic Earth): chalk pillars of Monument Rocks, Kansas, James St.~John, CC BY 2.0; the Bonarelli level, Picenus, CC BY-SA 4.0.
  \item Ch.~20 (fossils and palaeoenvironments): \emph{Waptia} and \emph{Marrella} from the Burgess Shale, James St.~John, CC BY 2.0; the Berlin specimen of \emph{Archaeopteryx}, H.~Raab, CC BY-SA 3.0; \emph{Zoophycos} at Sagres, Amada44, CC BY 4.0.
  \item Ch.~21 (atmospheric thermodynamics): anvil of a supercell thunderstorm, Stephen Corfidi, NOAA, public domain; radiosonde soundings of the Integrated Global Radiosonde Archive, I. Durre and others, NOAA NCEI, public domain.
  \item Ch.~22 (weather systems): a low-pressure system over Iceland, NASA/GSFC, MODIS Rapid Response Team, Jacques Descloitres, public domain; the eye of Hurricane Isabel from the International Space Station, Ed Lu, NASA, public domain.
  \item Ch.~23 (climate variability): North Atlantic Oscillation and Oceanic Ni\~no indices, NOAA Climate Prediction Center, public domain.
  \item Ch.~24 (climate archives and proxies): annual layers of the GISP2 ice core at 1855~m, Anthony Gow, U.S. Army Corps of Engineers CRREL (NOAA Paleoclimatology), public domain; cross-section of a travertine speleothem, Diamond Caverns, Kentucky, James St.~John, CC BY 2.0; the LR04 benthic stack of L.\,E. Lisiecki and M.\,E. Raymo (2005), the EPICA Dome C temperature record of J. Jouzel and others (2007) and the Antarctic \ce{CO2} composite of B. Bereiter and others (2015), NOAA NCEI Paleoclimatology, public domain.
  \item Ch.~25 (orbital forcing and the ice ages): orbital elements and insolation of Berger and Loutre (1991), NOAA NCEI Paleoclimatology, public domain.
  \item Ch.~26 (climate over geological time): Cenozoic benthic isotope record CENOGRID, T. Westerhold and others (2020), PANGAEA, CC BY 4.0; Phanerozoic \ce{CO2} compilation of G.\,L. Foster, D.\,L. Royer and D.\,J. Lunt (2017), \emph{Nature Communications}, CC BY 4.0.
  \item Ch.~27 (anthropogenic climate change): IPCC Sixth Assessment Report, Working Group~I, Annex~III extended data, CC BY 4.0.
  \item Ch.~28 (geological hazards and risk): debris-flow deposits at Caraballeda, Venezuela, Lawson Smith, U.S. Army Corps of Engineers, public domain; a mosque standing after the 2004 tsunami in Aceh, USGS, public domain.
  \item Ch.~29 (ore deposits): Bingham Canyon Mine from the International Space Station, NASA (Expedition~15 crew), public domain; gossan of the Flambeau Mine, Wisconsin, James St.~John, CC BY 2.0; grade--tonnage data of D.\,A. Singer, V.\,I. Berger and B.\,C. Moring (2008), USGS, public domain.
  \item Ch.~30 (energy geology): geothermal power plant near Geyserville, California, U.S. Department of Defense, public domain; carbon dioxide densities from the NIST Chemistry WebBook (Standard Reference Database 69).
  \item Ch.~31 (field and lab methods II): poikilitic augite enclosing plagioclase, Strekeisen, CC BY-SA 4.0; thin section of a garnet--mica schist, Jackdann88, CC BY 3.0; electron microprobe at the Geological Survey of Finland, Kallerna, CC BY-SA 4.0; silica results for rock reference materials, J.\,E. Taggart and D.\,F. Siems (2002), USGS, public domain.
\end{itemize}
